\section{How the results could form a quasi-polynomial portfolio}
\label{sec:portfolio}

\subsection{Three sufficient regimes}

The previous sections give different, explicitly checkable sufficient
conditions. They should be combined by a logical decision policy and a
complete cost account, rather than by adding their most favorable
exponents informally.

\begin{corollary}[Explicit sufficient regimes]
\label{cor:regimes}
For diagrams with $n$ crossings, any one of the following conditions
gives a deterministic $n^{O(\log n)}$ recognition method, including its
stated construction cost:
\begin{enumerate}[label=(\roman*)]
\item A verified presentation can be constructed in polynomial time with
      $r=O(\log n)$ and explicit relator length $N=n^{O(1)}$.
\item A verified word circuit of polynomial description length and a cut
      can be constructed in polynomial time with
      $(r+t)\log(d+2)=O((\log n)^2)$.
\item An order for the first-jet decision scan can be constructed in
      polynomial time whose actual one-object gap and frontier obey
      $g+w=O((\log n)^2)$, with all temporary allocation and the terminal
      segment included.
\end{enumerate}
No condition is asserted to hold on arbitrary diagrams. In (i) and (ii),
the decision method is the theoretical exact real procedure, not an
asymptotic guarantee for the supplied practical solver.
\end{corollary}

\begin{proof}
In (i), the logarithm of the nonpolynomial factor is
$O(r\log(2+N/r))=O((\log n)^2)$. Condition (ii) is precisely the
exponent of Theorem~\ref{thm:su2-hybrid}. Condition (iii) follows from
the scan bound in Section~\ref{sec:first-jet-pre-schur}. In each case
the additional polynomial factors are absorbed.
\end{proof}

The supplied-diagram condition $b(D)=O(\log n)$ is a concrete instance
of (i) whose construction is already implemented. Fully retaining a
word circuit with $r+s=O(\log^2 n)$ is a concrete instance of (ii).
Conditions (i) and (iii) describe different geometric and algebraic
phenomena: neither a small bridge invariant nor a small matrix response
may be substituted for the measured parameters in their proofs.

\subsection{Certified reductions and fair execution}

\begin{proposition}[A portfolio sufficient condition]
\label{prop:portfolio}
Suppose a polynomial-time verified reduction converts an input diagram
into polynomially many knot instances of total polynomial size, with the
original knot trivial if and only if every resulting knot is trivial.
For each instance suppose at least one member of a fixed finite collection
of sound algorithms terminates within $n^{O(\log n)}$ steps, with uniform
constants. Then a fair execution of that collection decides the original
input in $n^{O(\log n)}$ time. Every proposed certificate must be checked
within the corresponding cost, and the instance-size bound refers to the
actual encoded inputs.
\end{proposition}

\begin{proof}
Simulate the fixed number of algorithms on each instance in round-robin
instruction blocks. Stop that instance when one supplies a checked
verdict. If an algorithm takes $C$ steps, the simulation gives it those
steps with only a fixed-factor scheduling overhead in a standard
multi-workspace machine model; a conventional bit-machine simulation
adds at most polynomial or logarithmic overhead, harmless here. Sum
over polynomially many instances. Soundness and the verified reduction
justify the conjunction of their unknot verdicts. A single checked
knotted instance permits earlier rejection.
\end{proof}

This elementary proposition is not a new structural theorem. It locates
the missing proof: every residual input must be covered by at least one
bounded regime. Finite budgets, empirical dispatch thresholds, and a
collection of examples do not establish that covering statement. A
bounded practical solver that can return unknown is a sound portfolio
member, but does not by itself supply the terminating algorithm required
by the proposition.

For a scan making $q$ boundary observations over $a$ primes, the new
response contribution has the schematic arithmetic cost
\[
 O\bigl(a n(W+1)^2+a q(W+1)^3\bigr).
\]
Source geometry and index arithmetic have the separately stated costs
in Section~\ref{sec:streaming-response}. If $a$ or $q$ is large, that
must be paid. The number of relative objects and cobordism compositions
remains in the rest of the scan's cost. This is why the response theorem
is an ingredient within a recognition algorithm, not a complete
recognizer to place by itself in the portfolio minimum.

\subsection{Correct integration contracts}

The delivered additions expose narrow interfaces with explicit verdict
semantics. A generic group presentation produces a representation query;
only a formula derived from a validated knot diagram is promoted to a
knot verdict. Defining-generator elimination preserves the surviving
positive meridians, which permits the shared scalar coordinate and the
meridional inequality. Arbitrary future Whitehead transformations must
use the general commutator test unless conjugacy provenance is retained.

The first-jet observer accepts whole components of the actual differential
graph. It does not treat a selected submatrix as a direct summand. Local
exhaustion discards uninstalled proposals, disables observation, and resumes
ordinary elimination; global exhaustion propagates. Replaced complexes
support total-rank decisions only. A homology-reporting API must continue
to compute the original graded complex.

The response module returns a field-specific cofactor response. Its
absorbing zero is proved for every admissible continuation, but only in
that field. A zero residue is not an integer zero certificate, and
compatibility with an unknot determinant is not a positive recognition
certificate. Completion geometry, signs, loop phases, and the anchor
remain part of the exact query.

\begin{center}\small
\begin{tabularx}{\textwidth}{@{}p{0.25\textwidth}YY@{}}
\toprule
New interface & Definite output means & Inconclusive or excluded behavior\\
\midrule
\path{su2.solve_formula} & Exact nonabelian representation exists, or exact infeasibility for its input formula & Solver unknown, timeout, missing solver, or compilation allowance.\\
\path{su2.su2_decide} & Knot verdict derived from validated diagram provenance & Practical solver trust is explicit; generic external presentations are not accepted as diagram provenance.\\
\texttt{first\_jet\_profile} & Two exact ranks of a whole one-matching component & No inference from first-jet identities alone that an arbitrary polynomial matrix squares to zero.\\
\path{first_jet_khovanov_decide} & Total-rank decision preserved through its replacements & Original bigraded homology is outside the API contract.\\
\path{compile_suffix_responses} & Every prepared suffix has an immutable exact response modulo the supplied prime & No automatic driver dispatch or general recognition-time claim.\\
\bottomrule
\end{tabularx}\end{center}

The default backend selection is unchanged. The new modules are optional
research interfaces because the measurements do not yet justify a uniform
production dispatch policy. One unrelated existing subprocess polling
defect is corrected by integrating a previously developed fix; its
provenance and regression test are documented in Appendix~\ref{sec:reproduce}.
